\documentclass[final,a4paper,twocolumn,10pt,polish]{article}
\usepackage{amsmath}
\usepackage{mathtools}
\usepackage{fontspec}
\usepackage{unicode-math}
% \setmainfont{STIX Two Text}
% \setmathfont{STIX Two Math}
\setmonofont{JetBrains Mono}[Scale=MatchLowercase]
\usepackage{braket}
\usepackage{babel}
\usepackage{luavlna}
\singlechars{polish}{AaIiOoUuWwZz}
\usepackage{microtype}
\usepackage{csquotes}
\DeclareQuoteAlias[guillemets*]{polish}{polish}
\usepackage{indentfirst}
\usepackage{enumitem}
\usepackage{comment}
% a4 = 210mm x 297mm
\usepackage[a4paper,textwidth=160mm,textheight=235mm,centering]{geometry}
\linespread{1.05}
\usepackage{graphicx}
\usepackage{float}
\usepackage{caption}
\usepackage{subcaption}
\captionsetup[figure]{name=Wykres}
\usepackage{tabularx}
\usepackage{booktabs}
\usepackage[nospace,polish]{varioref}
\usepackage{xcolor}
\usepackage{hyperref}
\usepackage{xurl}
\hypersetup{
    colorlinks=true,
    linkcolor=black,
    urlcolor=red,
    bookmarksnumbered=true
}
\usepackage{siunitx}
\DeclareSIUnit{\litre}{l}
\DeclareSIUnit{\arbitrary}{a.u.}
\DeclareSIUnit{\angstrom}{\AA}
\sisetup{
    per-mode=symbol,
    exponent-product=\times,
    range-units=repeat,
    range-open-phrase={od },
    output-decimal-marker={.}
}
\AtBeginDocument{
    \let\square\mdlgwhtsquare
    \let\blacksquare\mdlgblksquare
    \let\varpropto\propto
}
\DeclarePairedDelimiter\abs{\lvert}{\rvert}
\usepackage{yfonts}

%\usepackage[version=4]{mhchem}
    
% \usepackage{tikz}
% \usetikzlibrary{positioning, arrows.meta, calc}

% \usepackage{listings}
% \lstset{
%   basicstyle=\ttfamily\small,
%   breaklines=false,
%   frame=single
% }

\title{\Huge{\textgoth{Sprawozdanie\\Lab 2}}}
\author{\huge{\textfrak{Maciej Flaga}}}
\date{
    \textgoth{Pazdziernik 2026}
}

\begin{document}
 
\maketitle
 
% \tableofcontents

\section{Energia}

\subsection{Metoda}

Zdefiniowano macierze zgodnie ze wzorami w treści zadania.

Rozwiązano uogólniony problem własny $\hat{H}\mathbf{c}_k=E_k\hat{S}\mathbf{c}_k$ wykorzystując do tego bibliotekę \texttt{GSL} i jej procedurę \texttt{gsl\_eigen\_gensymmv}.
Otrzymano w ten sposób $N$ energii własnych oraz $N$ zestawów wektorów własnych będących współczynnikami rozwinięcia funkcji falowych w dobranej bazie gaussów.

\subsection{Wyniki}

Pierwsze 5 poziomów energetycznych w funkcji $d$ i $N$ przedstawiają wykresy od \ref{Ek0} do \ref{Ek4}.

\begin{figure}
    \centering
    \includegraphics[width=\linewidth]{Ek0.png}
    \caption{Pierwszy poziom energetyczny w funkcji $d$ i $N$.}
    \label{Ek0}
\end{figure}

\begin{figure}
    \centering
    \includegraphics[width=\linewidth]{Ek1.png}
    \caption{Drugi poziom energetyczny w funkcji $d$ i $N$.}
    \label{Ek1}
\end{figure}

\begin{figure}
    \centering
    \includegraphics[width=\linewidth]{Ek2.png}
    \caption{Trzeci poziom energetyczny w funkcji $d$ i $N$.}
    \label{Ek2}
\end{figure}

\begin{figure}
    \centering
    \includegraphics[width=\linewidth]{Ek3.png}
    \caption{Czwarty poziom energetyczny w funkcji $d$ i $N$.}
    \label{Ek3}
\end{figure}

\begin{figure}
    \centering
    \includegraphics[width=\linewidth]{Ek4.png}
    \caption{Piąty poziom energetyczny w funkcji $d$ i $N$.}
    \label{Ek4}
\end{figure}

Zbadano również błąd energii podstawowej oraz błąd sumy pięciu pierwszych energii w funkcji $d$ i $N$.
Wyniki przedstawione na wykresach \ref{dE0} i \ref{dEsum}.

\begin{figure}
    \centering
    \includegraphics[width=\linewidth]{pkt1dE0.png}
    \caption{Błąd energii podstawowej.}
    \label{dE0}
\end{figure}

\begin{figure}
    \centering
    \includegraphics[width=\linewidth]{pkt1dEsum.png}
    \caption{Błąd sumy energii pierwszych pięciu poziomów.}
    \label{dEsum}
\end{figure}

\subsection{Wnioski}

Widzimy, że im gęstsza baza gaussów tym dokładniejszy wynik energii.

\section{Uwarunkowanie bazy}

\subsection{Metoda}

Obliczono $\kappa=s_{max}/s_{min}$, gdzie $\mathbf{s}$ jest zbiorem wartości własnych macierzy przekrywania $\hat{S}$, dla kilku wartości $N$, $d$ oraz $\beta$.

\subsection{Wyniki}

Wyniki przedstawiają wykresy od \ref{kb0} do \ref{kb4}.

\begin{figure}
    \centering
    \includegraphics[width=\linewidth]{kb0.png}
    \caption{$\kappa(d)$; $\beta=0.1$}
    \label{kb0}
\end{figure}

\begin{figure}
    \centering
    \includegraphics[width=\linewidth]{kb1.png}
    \caption{$\kappa(d)$; $\beta=0.2$}
    \label{kb1}
\end{figure}

\begin{figure}
    \centering
    \includegraphics[width=\linewidth]{kb2.png}
    \caption{$\kappa(d)$; $\beta=0.3$}
    \label{kb2}
\end{figure}

\begin{figure}
    \centering
    \includegraphics[width=\linewidth]{kb3.png}
    \caption{$\kappa(d)$; $\beta=0.4$}
    \label{kb3}
\end{figure}

\begin{figure}
    \centering
    \includegraphics[width=\linewidth]{kb4.png}
    \caption{$\kappa(d)$; $\beta=0.5$}
    \label{kb4}
\end{figure}

\subsection{Wnioski}

Widzimy, że w większości przypadków $\kappa$ nie zależy od $d$.
Ma to sens ponieważ rozpisując $\alpha$ i $X_i$ w definicji macierzy $\hat{S}$ dostajemy:

\begin{align*}
    S_{ij} &= \sqrt{\frac{\pi}{2\alpha}}\exp{\left[ -\frac{\alpha^2}{2\alpha}(X_i-X_j)^2 \right]} = \\
    &=\sqrt{\frac{\pi}{2\frac{\beta}{d^2}}}\exp{\left[ -\frac{\beta}{2d^2}d^2(i-j)^2 \right]} = \\
    &=\sqrt{\frac{\pi d^2}{2\beta}}\exp{\left[ -\frac{\beta}{2}(i-j)^2 \right]}
\end{align*}

Wartość $d$ jest jedynie współczynnikiem skalującym, który skraca się podczas dzielenia $s_{max}$ przez $s_{min}$.

Zbadano zależność $\kappa(\beta)$ co przedstawia wykres \ref{kappa}.

\begin{figure}
    \centering
    \includegraphics[width=\linewidth]{kappa.png}
    \caption{$\kappa(\beta)$}
    \label{kappa}
\end{figure}

Dla niskich wartości $\beta$ obserwujemy bardzo złe uwarunkowanie ponieważ funkcje bazowe nachodzą na siebie i baza nie jest niezależna.
Skutkuje to również niestabilnością numeryczną co można zauważyć na wykresach \ref{kb0} i \ref{kappa}.

Na wykresach \ref{dE0} i \ref{dEsum} obserwujemy również zwiększenie błędu dla największej wartości $d$. 
Jest to spowodowane tym, że odległe funkcje gaussa mają słabe pokrycie między sobą i nie są w stanie dobrze oddać kształtu rozwiązania.

\section{Funkcje falowe}

\subsection{Metoda}

Zbadano też samo rozwiązanie problemu oscylatora harmonicznego dla $N=31$ oraz $d=\qty{4}{\nano\meter}$.

\subsection{Wyniki}

Wykres \ref{psi} przedstawia pięć najniższych funkcji falowych na dziedzinie $x\in[-70,70]\ \unit{\nano\meter}$.

\begin{figure}
    \centering
    \includegraphics[width=\linewidth]{psi.png}
    \caption{Pięć najniższych stanów oscylatora harmonicznego.}
    \label{psi}
\end{figure}

\subsection{Wnioski}

Rozwiązanie wydaje się być zgodne z oczekiwanymi wygaszanymi wielomianami Hermite'a.
Kolejne rozwiązania charakteryzują się:

\begin{itemize}
    \item przemienną parzystością
    \item inkrementującą liczbą węzłów
    \item zanikiem na brzegach
\end{itemize}

Program przy odpaleniu również wypisuje do terminala wartości $\bra{\mathbf{c}_0}\hat{S}\ket{\mathbf{c}_0}$ oraz $\bra{\mathbf{c}_0}\hat{S}\ket{\mathbf{c}_{N/2}}$, które za każdym razem wynoszą lub są blisko odpowiednio $1$ i $0$.

\end{document}
